\documentclass[11pt]{article}

\usepackage[margin=1in]{geometry}
\usepackage{amsmath,amssymb,amsthm}
\usepackage{booktabs}
\usepackage{array}
\usepackage{xcolor}
\usepackage{listings}
\usepackage[colorlinks=true,linkcolor=blue!55!black,citecolor=blue!55!black,urlcolor=blue!55!black]{hyperref}

\emergencystretch=3em  % relax justification to avoid stubborn overfull lines in text

% ---- listing style (generic, no language-specific parser needed) ----
\definecolor{codebg}{gray}{0.97}
\lstset{
  basicstyle=\ttfamily\footnotesize,
  backgroundcolor=\color{codebg},
  breaklines=true,
  columns=fullflexible,
  keepspaces=true,
  frame=single,
  framesep=4pt,
  rulecolor=\color{gray!40},
  showstringspaces=false,
  literate={≠}{{$\neq$}}1 {⁻}{{$^-$}}1 {Ĝ}{{$\hat G$}}1 {σ}{{$\sigma$}}1 {θ}{{$\theta$}}1 {⊥}{{$\perp$}}1 {⊤}{{$\top$}}1
}

\theoremstyle{plain}
\newtheorem{proposition}{Proposition}
\newtheorem{lemma}{Lemma}
\newtheorem{corollary}{Corollary}
\theoremstyle{definition}
\newtheorem{definition}{Definition}
\newtheorem{remark}{Remark}

\newcommand{\R}{\mathbb{R}}
\newcommand{\tr}{\operatorname{tr}}
\newcommand{\diag}{\operatorname{diag}}
\newcommand{\Ghat}{\hat{G}}
\newcommand{\had}{\odot} % Hadamard product

\title{\textbf{Acceleration of Network Metrics in \texttt{MaxEntropyGraphs.jl}}\\[2pt]
\large Algorithmic reformulations, mathematical equivalence, autodiff and memory considerations}
\author{Implementation note}
\date{\today}

\begin{document}
\maketitle

\begin{abstract}
We document the acceleration of the graph-metric kernels in \texttt{src/metrics.jl} of
\texttt{MaxEntropyGraphs.jl}. The metrics are evaluated both \emph{analytically} on a dense expected
(bi)adjacency probability matrix $\Ghat$ and \emph{empirically} on many sampled graphs. We reformulate the
average nearest-neighbour degree (ANND), the triangle count, the thirteen directed three-node motifs, the
count of ``pure'' squares, and the bipartite $V$-motif count. For each we give the exact mathematical
identity underlying the reformulation, a proof of computational equivalence to the original loop
(valid for \emph{real-valued} matrices, not merely $0/1$ matrices), the resulting complexity and
peak-memory behaviour, and the constraints imposed by the two cross-cutting requirements of the package:
(i) the kernels lie on the reverse-/forward-mode automatic-differentiation (AD) path used to compute metric
variances, and (ii) they are frequently mapped over many graphs inside an already thread-parallel outer
loop, so per-call memory must not scale with the thread count. We close with the validation methodology
(bit-exact integer counts, \texttt{isapprox} float expectations, finite-difference gradient checks) and
empirical timing/memory results.
\end{abstract}

\tableofcontents

% =====================================================================================================
\section{Context and constraints}
\label{sec:context}

\texttt{MaxEntropyGraphs.jl} fits maximum-entropy null models (UBCM, DBCM, BiCM) to networks and uses them
to assess the statistical significance of topological quantities. A fitted model exposes a dense
\emph{expected adjacency matrix} $\Ghat\in\R^{N\times N}$ (for the bipartite BiCM, an expected biadjacency
matrix $\Ghat\in\R^{n_\perp\times n_\top}$), whose entries are edge probabilities in $[0,1]$. For UBCM/DBCM
$\Ghat$ is real-valued with a \emph{zero diagonal}; for UBCM it is additionally symmetric.

There are two ways a metric $X$ is consumed:
\begin{enumerate}
  \item \textbf{Analytical route.} Evaluate $X$ directly on $\Ghat$, giving $X(\langle G\rangle)$, an
  approximation of $\mathbb{E}[X]$. The associated variance is obtained by the delta method,
  $\sigma_X = \sqrt{\sum_{ij}(\sigma_{ij}\,\partial X/\partial \Ghat_{ij})^2}$, where the gradient
  $\nabla_{\Ghat} X$ is computed by AD (default ReverseDiff, also ForwardDiff and Zygote).
  \item \textbf{Sampling route.} Draw many realisations $G_s\sim$ model and evaluate $X(G_s)$ on each,
  then form Monte-Carlo means, standard deviations and $z$-scores.
\end{enumerate}

Three properties of this usage shape every reformulation below.

\paragraph{(C1) Parallelism is at the outer level.} The batch samplers \texttt{rand(m,n)} and the
significance sweeps map a metric over $n$ graphs with \texttt{Threads.@threads}. Introducing inner
threading in a metric would nest inside an already-saturated pool. \emph{We therefore add no inner
threading}; speed-ups come from lower asymptotic complexity, from BLAS, and from removing branch overhead.

\paragraph{(C2) Memory must not scale with the thread count.} If a kernel materialises an $N\times N$
intermediate, then under $T$ threads the transient footprint is $T\cdot O(N^2)$. For $N=10^4$ a single dense
$N\times N$ \texttt{Float64} matrix is already $800$\,MB. \emph{We prefer closed forms and column-streaming
formulations whose \emph{peak} extra memory is $O(N)$}, reserving materialised products for cases where they
are unavoidable and then bounding and reusing them.

\paragraph{(C3) AD-safety.} Because $X$ is differentiated with respect to $\Ghat$, the matrix methods must
remain differentiable. ForwardDiff and ReverseDiff substitute a custom element type
(\texttt{Dual}, \texttt{TrackedReal}) and re-dispatch; Zygote is source-to-source and differentiates the
method chosen for the \emph{concrete} \texttt{Float64} matrix. In particular, in-place mutation
(\texttt{mul!}) breaks Zygote even though ReverseDiff/ForwardDiff tolerate it, because the latter never
reach the mutating concrete method. All reformulations below are written non-mutating on the paths Zygote
can reach.

\begin{remark}[Design pattern]
Each metric public method performs input checks and then dispatches to an internal kernel. Where a
memory-frugal concrete-float variant is worthwhile, it is a specialisation on
\texttt{AbstractMatrix\{<:BLAS.BlasFloat\}}; \texttt{Dual}/\texttt{TrackedReal} element types are $<:$\,\texttt{Real}
but \emph{not} $<:$\,\texttt{BlasFloat}, so AD inputs fall through to the generic (allocating, non-mutating)
kernel automatically. The model methods (\texttt{triangles(m::UBCM)}, etc.) merely forward to the matrix
method on $\Ghat$ and are unchanged.
\end{remark}

% =====================================================================================================
\section{Notation}
\label{sec:notation}
$A\in\R^{N\times N}$ denotes an adjacency or expected-adjacency matrix with $A_{ii}=0$. $J$ is the all-ones
matrix, $\had$ the Hadamard (entrywise) product, $\|\cdot\|_F$ the Frobenius norm, $\mathbf 1$ the all-ones
vector. Row/column sums (degrees) are $k^{\text{col}}_i=\sum_j A_{ji}=(\mathbf 1^\top A)_i$ and
$k^{\text{row}}_i=\sum_j A_{ij}=(A\mathbf 1)_i$; for symmetric $A$ they coincide and we write $k_i$. For a
biadjacency matrix $B\in\R^{n_\perp\times n_\top}$ we write column sums $s=B^\top\mathbf 1$ and row sums
$r=B\mathbf 1$.

% =====================================================================================================
\section{ANND: from $O(N^3)$ to $O(N^2)$}
\label{sec:annd}

\begin{definition}
For a node $i$ with degree $k_i>0$, the average nearest-neighbour degree is
$\mathrm{ANND}_i(A)=\frac{1}{k_i}\sum_{j} A_{ij}\,k_j$, and $\mathrm{ANND}_i=0$ if $k_i=0$.
\end{definition}

The original vector method called a per-node routine that recomputed \emph{all} column sums inside a
\texttt{mapreduce}; over $N$ nodes this is $O(N^3)$.

\begin{proposition}[Equivalence]\label{prop:annd}
Let $k=\mathbf 1^\top A$ (column sums). Then for every node $i$ with $k_i>0$,
\[
\mathrm{ANND}_i(A) \;=\; \frac{(A k)_i}{k_i}.
\]
\end{proposition}
\begin{proof}
$(Ak)_i=\sum_j A_{ij}k_j$ by definition of matrix--vector product, which is exactly the numerator of
$\mathrm{ANND}_i$; dividing by $k_i$ gives the claim. The in/out variants use $k=\mathbf 1^\top A$
(in-degrees) and $k=A\mathbf 1$ (out-degrees) respectively.
\end{proof}

Thus the whole vector is one matrix--vector product plus an elementwise divide with a zero-degree guard:
\begin{lstlisting}
k = vec(sum(A, dims=1))          # column sums (degrees), computed ONCE
num = A * k                      # single gemv, O(N^2) time, O(N) memory
return [iszero(k[i]) ? zero(...) : num[i]/k[i] for i in vs]
\end{lstlisting}
This lowers the cost from $O(N^3)$ to $O(N^2)$ and the extra memory to $O(N)$. Integer exactness is immediate:
for an integer/Boolean $A$, $Ak$ is exact integer arithmetic (order-independent), so the result is
bit-identical to the graph-based computation. On the AD path all degrees are strictly positive, so the guard
branch is never taken; the guard uses \emph{instance}-level \texttt{zero}/\texttt{oneunit}, avoiding
type-level calls that are undefined for some tracked element types.

% =====================================================================================================
\section{Triangles: $\tr(A^3)/6$ with $O(N)$ peak memory}
\label{sec:triangles}

\begin{definition}
The number of (expected) triangles is
$T(A)=\tfrac16\sum_{i,j,k\ \text{distinct}} A_{ij}A_{jk}A_{ki}$.
\end{definition}

\begin{proposition}[Equivalence]\label{prop:tri}
If $A_{ii}=0$ for all $i$, then $T(A)=\tfrac16\tr(A^3)$.
\end{proposition}
\begin{proof}
$\tr(A^3)=\sum_{i,j,k}A_{ij}A_{jk}A_{ki}$ sums over \emph{all} index triples. Consider a term with a
coincidence. If $i=j$ then the factor $A_{ij}=A_{ii}=0$; if $j=k$ then $A_{jk}=A_{jj}=0$; if $k=i$ then
$A_{ki}=A_{ii}=0$. Hence every term with a repeated index vanishes and
$\tr(A^3)=\sum_{i,j,k\ \text{distinct}}A_{ij}A_{jk}A_{ki}$. Dividing by $6$ gives $T(A)$.
\end{proof}

Note the proof uses \emph{only} the zero diagonal, so it holds for arbitrary real $A$ (in particular for
$\Ghat$), not merely for $0/1$ matrices. Two evaluation forms are used.

\paragraph{Generic / AD path.} For symmetric $A$, $\tr(A^3)=\sum_{i,k}A_{ik}(A^2)_{ik}=\langle A,A^2\rangle$
(Frobenius inner product), because $A_{ik}=A_{ki}$. Hence $T=\langle A,A A\rangle/6$, i.e.
\texttt{dot(A, A*A)/6}. This is a single BLAS-3 \texttt{gemm} for \texttt{Float64} and is differentiable for
\texttt{Dual}/\texttt{TrackedReal}; it is the form ForwardDiff/ReverseDiff/Zygote traverse. It allocates one
$N\times N$ temporary.

\paragraph{Memory-frugal concrete path (default for \texttt{BlasFloat}).} Streaming the diagonal of $A^3$
column by column,
\[
\tr(A^3)=\sum_i (A^3)_{ii}=\sum_i \big\langle A_{i,:}\,,\ A\,A_{:,i}\big\rangle ,
\]
which uses one \texttt{gemv} per column and a single length-$N$ vector, so the \emph{peak} extra memory is
$O(N)$ instead of $O(N^2)$. This is the choice that keeps the batch/threaded workload (C2) from multiplying
an $N\times N$ temporary by the thread count. It is written with non-mutating \texttt{A*A[:,i]} (not
\texttt{mul!}) so Zygote --- which differentiates the concrete-\texttt{Float64} method --- remains valid; the
row/column split $\langle A_{i,:},AA_{:,i}\rangle=(A^3)_{ii}$ is correct without symmetry.

\begin{remark}[Speed/memory trade-off]
The streaming form is BLAS-2 and yields a $\sim$5--6$\times$ speed-up over the original branchy triple loop
(Table~\ref{tab:bench}). The generic \texttt{dot(A,A*A)} form is BLAS-3 and faster still, at the cost of an
$O(N^2)$ peak temporary. We default to the memory-frugal form; switching is a one-line change.
\end{remark}

\begin{corollary}[Gradient sanity]
$\nabla_A \tr(A^3)=3(A^2)^\top$, hence $\nabla_A T=\tfrac12 (A^2)^\top=\tfrac12 A^2$ for symmetric $A$. This
analytic gradient is reproduced by all three AD backends (validated by finite differences,
\S\ref{sec:validation}).
\end{corollary}

% =====================================================================================================
\section{Directed three-node motifs: an inclusion--exclusion matrix form}
\label{sec:motifs}

Each of the thirteen connected directed three-node motifs is a sum over ordered distinct triples of a product
of three \emph{dyadic indicators} of the pair states ``$i\!\to\!j$ only'', ``$j\!\to\!i$ only'',
``reciprocated'', ``absent'':
\begin{gather*}
a_{\to}(i,j)=A_{ij}(1-A_{ji}),\qquad
a_{\leftarrow}(i,j)=(1-A_{ij})A_{ji},\\
a_{\leftrightarrow}(i,j)=A_{ij}A_{ji},\qquad
a_{\varnothing}(i,j)=(1-A_{ij})(1-A_{ji}).
\end{gather*}
A motif with dyad functions $(f_1,f_2,f_3)$ is
\begin{equation}\label{eq:motif}
M=\sum_{i,j,k\ \text{distinct}} f_1(i,j)\,f_2(j,k)\,f_3(k,i).
\end{equation}
Here $J$ is the all-ones matrix; recall $A_{ii}=0$.
The original implementation evaluated \eqref{eq:motif} with a branchy $O(N^3)$ triple loop, once per motif.

\paragraph{Indicator matrices.} Define, entrywise over all $(i,j)$,
\begin{gather*}
P=A\had(J-A^\top)\ (a_{\to}),\qquad
Q=P^\top\ (a_{\leftarrow}),\\
R=A\had A^\top\ (a_{\leftrightarrow}),\qquad
Z=(J-A)\had(J-A^\top)\ (a_{\varnothing}).
\end{gather*}
Because $A_{ii}=0$, the diagonals satisfy $\diag(P)=\diag(Q)=\diag(R)=0$ and $\diag(Z)=\mathbf 1$
(since $(1-A_{ii})^2=1$). These four matrices are built once from $A$ and $A^\top$ (with $Q=P^\top$ a lazy
transpose, no allocation) and reused across all thirteen motifs.

\begin{proposition}[Distinct-index inclusion--exclusion]\label{prop:motif}
For matrices $F_1,F_2,F_3\in\R^{N\times N}$,
\begin{multline*}
\sum_{i,j,k\ \text{distinct}} (F_1)_{ij}(F_2)_{jk}(F_3)_{ki}
=\tr(F_1F_2F_3)
-\sum_a (F_1)_{aa}(F_2F_3)_{aa}\\
-\sum_a (F_2)_{aa}(F_3F_1)_{aa}
-\sum_a (F_3)_{aa}(F_1F_2)_{aa}
+2\sum_a (F_1)_{aa}(F_2)_{aa}(F_3)_{aa}.
\end{multline*}
\end{proposition}
\begin{proof}
Write $T=\sum_{i,j,k}(F_1)_{ij}(F_2)_{jk}(F_3)_{ki}=\tr(F_1F_2F_3)$ (the unrestricted sum). The distinct-index
sum removes the union of the three coincidence events $E_{ij}\!:i=j$, $E_{jk}\!:j=k$, $E_{ki}\!:k=i$. By
inclusion--exclusion,
$\Sigma_{\text{distinct}}=T-\big(\Sigma_{E_{ij}}+\Sigma_{E_{jk}}+\Sigma_{E_{ki}}\big)+\big(\text{pairwise
intersections}\big)-\Sigma_{E_{ij}\cap E_{jk}\cap E_{ki}}$. Any two of the events force $i=j=k$, so each
pairwise intersection equals $\Sigma_{i=j=k}$ and so does the triple intersection; there are three pairwise
terms, giving $+3\Sigma_{iii}-\Sigma_{iii}=+2\Sigma_{iii}$. Finally,
$\Sigma_{E_{ij}}=\sum_{i,k}(F_1)_{ii}(F_2)_{ik}(F_3)_{ki}=\sum_a (F_1)_{aa}(F_2F_3)_{aa}$, and analogously for
the other two events and the triple coincidence.
\end{proof}

\begin{corollary}\label{cor:motif}
Since only $Z$ has a nonzero diagonal (all ones) and no motif contains two $a_\varnothing$ factors, the triple
term vanishes and at most one correction survives, equal to a single trace of the two \emph{other} factors:
\[
M=\tr(F_1F_2F_3)-
\begin{cases}
\tr(F_2F_3) & \text{if } F_1=Z,\\
\tr(F_3F_1) & \text{if } F_2=Z,\\
\tr(F_1F_2) & \text{if } F_3=Z,\\
0 & \text{if no factor is } Z.
\end{cases}
\]
\end{corollary}

Using $\tr(XY)=\langle X,Y^\top\rangle$ and $\tr(XYW)=\langle XY,W^\top\rangle$, together with
$Q=P^\top$, $R^\top=R$, $Z^\top=Z$, each motif reduces to one \texttt{gemm} and inner products.
Table~\ref{tab:motifs} lists all thirteen (matching the dyad triples hard-coded in the source).

\begin{table}[t]
\centering
\small
\begin{tabular}{clll}
\toprule
$M$ & dyads $(f_1,f_2,f_3)$ & matrices $(F_1,F_2,F_3)$ & closed form \\
\midrule
$M_{1}$  & $(a_\leftarrow,a_\to,a_\varnothing)$        & $(Q,P,Z)$ & $\langle QP,Z\rangle-\langle Q,Q\rangle$ \\
$M_{2}$  & $(a_\leftarrow,a_\leftarrow,a_\varnothing)$ & $(Q,Q,Z)$ & $\langle QQ,Z\rangle-\langle Q,P\rangle$ \\
$M_{3}$  & $(a_\leftarrow,a_\leftrightarrow,a_\varnothing)$ & $(Q,R,Z)$ & $\langle QR,Z\rangle-\langle Q,R\rangle$ \\
$M_{4}$  & $(a_\leftarrow,a_\varnothing,a_\to)$        & $(Q,Z,P)$ & $\langle QZ,Q\rangle-\langle P,P\rangle$ \\
$M_{5}$  & $(a_\leftarrow,a_\to,a_\to)$                & $(Q,P,P)$ & $\langle QP,P^\top\rangle$ \\
$M_{6}$  & $(a_\leftarrow,a_\leftrightarrow,a_\to)$    & $(Q,R,P)$ & $\langle QR,P^\top\rangle$ \\
$M_{7}$  & $(a_\to,a_\leftrightarrow,a_\varnothing)$   & $(P,R,Z)$ & $\langle PR,Z\rangle-\langle P,R\rangle$ \\
$M_{8}$  & $(a_\leftrightarrow,a_\leftrightarrow,a_\varnothing)$ & $(R,R,Z)$ & $\langle RR,Z\rangle-\langle R,R\rangle$ \\
$M_{9}$  & $(a_\to,a_\to,a_\to)$                       & $(P,P,P)$ & $\langle PP,P^\top\rangle=\tr(P^3)$ \\
$M_{10}$ & $(a_\leftrightarrow,a_\to,a_\to)$           & $(R,P,P)$ & $\langle RP,P^\top\rangle$ \\
$M_{11}$ & $(a_\leftrightarrow,a_\leftarrow,a_\to)$    & $(R,Q,P)$ & $\langle RQ,P^\top\rangle$ \\
$M_{12}$ & $(a_\leftrightarrow,a_\leftrightarrow,a_\to)$ & $(R,R,P)$ & $\langle RR,P^\top\rangle$ \\
$M_{13}$ & $(a_\leftrightarrow,a_\leftrightarrow,a_\leftrightarrow)$ & $(R,R,R)$ & $\langle RR,R\rangle=\tr(R^3)$ \\
\bottomrule
\end{tabular}
\caption{The thirteen directed motifs and their closed forms
($\langle X,Y\rangle=\sum_{ab}X_{ab}Y_{ab}$).
$M_{13}=\tr(R^3)$ counts each reciprocated triangle $6$ times, matching the graph-based convention.}
\label{tab:motifs}
\end{table}

Each motif costs one $O(N^3)$ \texttt{gemm} instead of a branchy $O(N^3)$ loop; the batched \texttt{motifs(A)}
builds $P,R,Z$ once and returns the whole spectrum. Peak memory is bounded by a small constant number of
$N\times N$ matrices (the base matrices plus one product), which is $O(N^2)$ but does \emph{not} grow with
the number of motifs.

\begin{remark}[Sparse inputs]
$J-A^\top$ densifies a sparse $A$, so the matrix form targets the dense $\Ghat$. Sparse realised graphs use
the neighbour-based graph methods $M_k(G)$, which are retained unchanged.
\end{remark}

% =====================================================================================================
\section{Squares: an $8\times$ symmetry reduction and why $o(N^4)$ is out of reach}
\label{sec:squares}

\begin{definition}
The number of ``pure'' squares (chordless $4$-cycles) is
\begin{equation}\label{eq:squares}
S(A)=\frac18\sum_{i,j,k,l\ \text{distinct}} A_{ij}A_{jk}A_{kl}A_{li}\,(1-A_{ik})(1-A_{lj}),
\end{equation}
i.e.\ a $4$-cycle $i\!-\!j\!-\!k\!-\!l\!-\!i$ with both diagonals $\{i,k\}$ and $\{j,l\}$ absent.
\end{definition}

\subsection{Dihedral symmetry and exact orbit reduction}

Write the summand of \eqref{eq:squares} as $g(i,j,k,l)=A_{ij}A_{jk}A_{kl}A_{li}(1-A_{ik})(1-A_{lj})$.

\begin{lemma}[Invariance]\label{lem:dihedral}
For symmetric $A$, $g$ is invariant under the dihedral group $D_4$ (order $8$) acting on the positions of the
labelled $4$-cycle, generated by the rotation $\rho:(i,j,k,l)\mapsto(j,k,l,i)$ and the reflection
$\tau:(i,j,k,l)\mapsto(i,l,k,j)$.
\end{lemma}
\begin{proof}
The backbone $\{A_{ij},A_{jk},A_{kl},A_{li}\}$ is the edge set of the cycle and is preserved (as a multiset)
by both generators. Under $\rho$ the chord factors become $(1-A_{jl})(1-A_{ik})$; under $\tau$ they become
$(1-A_{ik})(1-A_{jl})$; using $A_{lj}=A_{jl}$ (symmetry) these equal $(1-A_{ik})(1-A_{lj})$. Hence $g$ is
unchanged.
\end{proof}

\begin{proposition}[Exact orbit sum]\label{prop:squares}
For distinct indices the $D_4$-orbit of $(i,j,k,l)$ has exactly $8$ elements, and
\[
S(A)=\sum_{\substack{i<j<l,\ k>i\\ k\notin\{j,l\}}} A_{ij}A_{jk}A_{kl}A_{li}\,(1-A_{ik})(1-A_{lj}).
\]
That is, summing one canonical representative per orbit and dropping the $1/8$ yields the exact value.
\end{proposition}
\begin{proof}
On distinct labels no non-identity element of $D_4$ fixes a tuple, so orbits have size $8$; by
Lemma~\ref{lem:dihedral} $g$ is constant on each orbit, hence $\sum_{\text{distinct}}g=8\sum_{\text{orbits}}g$
and $S=\tfrac18\sum_{\text{distinct}}g=\sum_{\text{orbits}}g$. It remains to see that the index constraints
select one representative per orbit. Fix the orbit and let $i$ be the minimum of the four labels; the two
cycle-neighbours of $i$ are the entries in positions $2,4$ and the diagonally opposite label $k$ is in
position $3$. Exactly two orbit elements place the minimum in position $1$ --- they are related by $\tau$
(which swaps positions $2,4$) --- and requiring the position-$2$ label to be smaller than the position-$4$
label ($j<l$) selects one. For a fixed $4$-set with minimum $i$, the three orbits correspond to the three
choices of the opposite label $k\in\{$the other three$\}$; the loop ranges $k>i$, $k\notin\{j,l\}$ over
exactly those choices. Thus the summation visits each orbit's canonical representative once.
\end{proof}

The kernel therefore performs $\approx N^4/8$ iterations, hoists the $k$-independent factors
$A_{ij}A_{li}(1-A_{lj})$ out of the inner loop, and accumulates into a scalar (no allocation):
\begin{lstlisting}
res = zero(eltype(A)); o = one(eltype(A))
for i in 1:n, j in i+1:n
    Aij = A[i,j]
    for l in j+1:n
        base = Aij * A[l,i] * (o - A[l,j])   # invariant over k
        iszero(base) && continue             # prunes only EXACT zeros (0/1 inputs)
        for k in i+1:n
            (k == j || k == l) && continue
            res += base * A[j,k] * A[k,l] * (o - A[i,k])
        end
    end
end
return res                                   # note: no /8
\end{lstlisting}
This is exact for real $A$ (it is Proposition~\ref{prop:squares} evaluated term by term), integer-exact for
$0/1$ inputs, and differentiable: the only control flow depends on indices or on \emph{exact} zeros, and on
the AD path $\Ghat\in(0,1)$ so \texttt{iszero(base)} is never true and no gradient contribution is dropped.

\subsection{Why there is no sub-$O(N^4)$ closed form}
Expanding $(1-A_{ik})(1-A_{lj})=1-A_{ik}-A_{lj}+A_{ik}A_{lj}$ gives $8S=S_0-2S_1+S_3$ (with $S_1=S_2$ by
symmetry), where $S_0$ counts $4$-cycles (backbone only), $S_1$ adds one chord, and
\[
S_3=\sum_{\text{distinct}} A_{ij}A_{jk}A_{kl}A_{li}A_{ik}A_{lj}=\hom(K_4,A)=24\cdot\#K_4 ,
\]
the homomorphism count of the \emph{complete} graph $K_4$ (all six pairs present). $S_0$ and $S_1$ admit
$O(N^3)$ closed forms in traces of powers of $A$ (e.g.\ $S_0=\|A^2\|_F^2-2\|p\|^2+\|A\had A\|_F^2$ with
$p_i=\sum_j A_{ij}^2$). $S_3$, however, is a treewidth-$3$ pattern: its evaluation does not reduce to a
polynomial in traces of powers of $A$ and is $\Omega(N^4)$ for dense $A$ by elementary means (fast
matrix-multiplication gives only $O(N^{\omega+1})$). Hence the pure-square count is intrinsically $\sim N^4$
for a dense $\Ghat$; the gain here is the exact $8\times$ constant-factor reduction plus branch elimination.

\subsection{Sparse fast path}
For a sparse $0/1$ adjacency matrix the dense kernel's scalar indexing is $O(\log)$ per access. We instead
dispatch \texttt{squares(A::SparseMatrixCSC\{<:Union\{Bool,Integer\}\})} to the neighbour-enumeration graph
algorithm (equivalent by the tested identity $\texttt{squares(G)}=\texttt{squares(adjacency\_matrix(G))}$),
which makes $N\gtrsim 10^4$ tractable where the dense $O(N^4)$ evaluation is not.

% =====================================================================================================
\section{Bipartite $V$-motifs: a marginal-sum closed form}
\label{sec:vmotifs}

\begin{definition}
For a biadjacency matrix $B$, the bottom-layer $V$-motif count is the number of shared-neighbour pairs
$V=\sum_{i<j}\langle B_{i,:},B_{j,:}\rangle=\sum_{i<j}(BB^\top)_{ij}$ (rows are bottom-layer nodes); the
top-layer count uses columns.
\end{definition}

The original method summed the pairwise dot products explicitly, an $O(N^2 M)$ computation reconstructing the
strict upper triangle of the Gram matrix.

\begin{proposition}[Equivalence]\label{prop:vmot}
With column sums $s=B^\top\mathbf 1$,
\[
V_\perp=\sum_{i<j}(BB^\top)_{ij}=\tfrac12\big(\|s\|^2-\|B\|_F^2\big),
\]
and the top-layer count is $\tfrac12(\|r\|^2-\|B\|_F^2)$ with row sums $r=B\mathbf 1$.
\end{proposition}
\begin{proof}
Let $\Sigma=\sum_{i,j}(BB^\top)_{ij}=\mathbf 1^\top BB^\top\mathbf 1=\|B^\top\mathbf 1\|^2=\|s\|^2$, and
$\tr(BB^\top)=\sum_i\|B_{i,:}\|^2=\|B\|_F^2$. Since $BB^\top$ is symmetric,
$\sum_{i<j}(BB^\top)_{ij}=\tfrac12(\Sigma-\tr(BB^\top))=\tfrac12(\|s\|^2-\|B\|_F^2)$.
\end{proof}

This never materialises $BB^\top$; it needs only the marginal sums and $\|B\|_F^2$, giving $O(NM)$ time and
$O(N+M)$ memory. Integer exactness follows because
$\|s\|^2-\|B\|_F^2=\sum_{i\neq j}(BB^\top)_{ij}=2\sum_{i<j}(BB^\top)_{ij}$ is even; the implementation uses
integer division $\div 2$ for integer element types (preserving the \texttt{Int} return type) and $/2$ for
floating/tracked types. The empirical speed-up is the largest of any metric here (Table~\ref{tab:bench}),
reaching $\sim$$2.5\times10^3$ at $N=10^3$, because both the asymptotic factor and the constant improve.

% =====================================================================================================
\section{Autodiff considerations}
\label{sec:ad}

The variance routine $\sigma_X$ differentiates $X$ with respect to $\Ghat$. The interaction with dispatch is
subtle and dictated the concrete forms above:
\begin{itemize}
  \item \textbf{ForwardDiff / ReverseDiff} replace \texttt{Float64} by \texttt{Dual}/\texttt{TrackedReal}.
  These are $<:$\,\texttt{Real} but not $<:$\,\texttt{BlasFloat}, so they reach the \emph{generic}
  non-mutating kernels (\texttt{dot(A,A*A)}, the matrix motif forms, the scalar squares loop, the ANND
  comprehension), all of which are differentiable.
  \item \textbf{Zygote} is source-to-source and differentiates the method selected for the \emph{concrete}
  \texttt{Matrix\{Float64\}}. This is why the memory-frugal triangles path must avoid \texttt{mul!}: an
  in-place write has no adjoint and aborts Zygote, whereas the non-mutating \texttt{A*A[:,i]} form
  differentiates cleanly. (ReverseDiff/ForwardDiff never reach this method, so a \texttt{mul!} version would
  have passed their tests while silently breaking Zygote --- exactly the regression the AD test-suite now
  guards against.)
\end{itemize}
All matrix metrics were checked to (i) run under \texttt{:ReverseDiff}, \texttt{:ForwardDiff} and
\texttt{:Zygote}, and (ii) match central finite differences (\S\ref{sec:validation}). For \texttt{squares},
whose evaluation is $O(N^4)$ over an $N^2$-dimensional input, reverse mode (a single tape pass) is the
recommended backend.

% =====================================================================================================
\section{Memory analysis}
\label{sec:memory}

Table~\ref{tab:memory} summarises the peak transient memory of each metric per call. The distinction between
\emph{peak} (simultaneously live) and \emph{churn} (total allocated, garbage-collected) matters under the
outer thread-parallel loop (C2): only peak is multiplied by the thread count. The triangles streaming form,
the ANND \texttt{gemv}, and the $V$-motif closed form all keep peak at $O(N)$; the motif spectrum keeps peak
at a bounded number of $N\times N$ matrices regardless of how many motifs are requested; the squares kernel
allocates nothing.

\begin{table}[t]
\centering
\small
\begin{tabular}{lccc}
\toprule
metric & original time & new time & new peak memory \\
\midrule
ANND                 & $O(N^3)$ & $O(N^2)$          & $O(N)$ \\
triangles (concrete) & $O(N^3)$ & $O(N^3)$ (BLAS-2) & $O(N)$ \\
triangles (generic/AD) & $O(N^3)$ & $O(N^3)$ (BLAS-3) & $O(N^2)$ \\
motifs (all 13)      & $13\cdot O(N^3)$ & $13\cdot O(N^3)$ (BLAS-3) & $O(N^2)$ (bounded, shared) \\
squares (dense)      & $O(N^4)$ & $O(N^4)/8$        & $O(1)$ \\
squares (sparse)     & $O(N^4)$ & neighbour enum.\  & --- \\
$V$-motifs           & $O(N^2M)$ & $O(NM)$          & $O(N+M)$ \\
\bottomrule
\end{tabular}
\caption{Complexity and peak extra memory. Peak (not churn) is what scales with the thread count.}
\label{tab:memory}
\end{table}

% =====================================================================================================
\section{Computational equivalence and validation}
\label{sec:validation}

Correctness is asserted at two tiers, mirroring how the package tests other quantities.

\paragraph{Integer counts --- bit-exact.} On $0/1$ graph and matrix inputs the accelerated kernels must equal
the graph-based counts and a reference (pre-acceleration) implementation \emph{exactly} (\texttt{==}). This
holds because all the reformulations reduce to integer arithmetic whose value is independent of summation
order (matrix products, inner products, the symmetric squares kernel). The test-suite checks
$\texttt{triangles}$, $\texttt{squares}$, all thirteen $M_k$, $\texttt{motifs}$ and $\texttt{V\_motifs}$ on the
built-in fixtures (karate/UBCM, \texttt{maspalomas}/DBCM, \texttt{corporateclub}/BiCM) and on
Barab\'asi--Albert / Erd\H{o}s--R\'enyi graphs, including the sparse squares fast path.

\paragraph{Float expectations --- \texttt{isapprox}.} On real matrices (random symmetric/directed matrices in
$[0,1]$ and the fitted $\Ghat$) the reformulations change the summation order, so they are validated against
the reference loops with a tight relative tolerance ($10^{-10}$). The affected documentation
\texttt{jldoctest} outputs were updated to rounded, platform-robust values.

\paragraph{Autodiff --- finite differences.} For fitted UBCM and DBCM models the reverse-mode gradient of
each metric with respect to $\Ghat$ is compared to a central finite-difference gradient at several entries
(relative tolerance $\sim10^{-4}$), and $\sigma_X$ is exercised end-to-end under all three backends.

The full package test-suite (759 tests) passes, including the new
\texttt{"metrics acceleration"} test-set.

% =====================================================================================================
\section{Empirical results}
\label{sec:empirical}

Table~\ref{tab:bench} reports head-to-head timings (Julia 1.12, single BLAS thread, to reflect the
batch/threaded-over-graphs usage) produced by \texttt{performance/metrics\_benchmarks.jl}. The
complexity-reducing metrics (ANND, $V$-motifs) show a speed-up that \emph{grows} with $N$ --- the signature of
a genuine asymptotic improvement --- while the constant-factor metrics (triangles, motifs, dense squares)
show stable multiplicative gains.

\begin{table}[t]
\centering
\footnotesize
\begin{tabular}{llrrr}
\toprule
metric & regime & $N$ & speed-up vs.\ original & new peak/behaviour \\
\midrule
ANND                & $O(N^3)\!\to\!O(N^2)$ & $100/400/1000$ & $25\times/82\times/192\times$ & $O(N)$; $0.14$\,s at $N{=}10^4$ \\
$V$-motifs          & $O(N^2M)\!\to\!O(NM)$ & $100/400/1000$ & $155\times/856\times/2482\times$ & $O(N{+}M)$; $0.026$\,s at $N{=}10^4$ \\
motifs (all 13)     & constant factor       & $50/150$        & $18\times/22\times$          & bounded shared $N\times N$ \\
squares (dense)     & constant factor       & $25/50/100/150$ & $\approx 12\text{--}13\times$ & $0$ extra bytes \\
triangles           & constant factor       & $100/400$       & $\approx 5\text{--}6\times$   & $O(N)$ peak \\
squares (sparse)    & new capability        & up to $5\!\times\!10^4$ & tractable ($11.6$\,s) & dense $O(N^4)$ infeasible \\
\bottomrule
\end{tabular}
\caption{Head-to-head timing (Julia 1.12, single-thread BLAS). The reusable harness supports a wider,
per-metric-adaptive size sweep and a second pass at the BLAS-library default thread count.}
\label{tab:bench}
\end{table}

% =====================================================================================================
\section{Incidental correctness fixes}
\label{sec:fixes}
Two latent issues surfaced during the work and were fixed in the same change:
\begin{itemize}
  \item \texttt{strength(G, i, T; dir)} previously computed and returned the \emph{entire} strength vector,
  ignoring the node argument \texttt{i}. It now returns node $i$'s scalar strength by indexing the
  appropriate marginal of the weight matrix.
  \item \texttt{ANND(m::UBCM)} re-ran an $O(N^2)$ \texttt{issymmetric} check on the by-construction symmetric
  $\Ghat$; the redundant check is now skipped.
\end{itemize}

% =====================================================================================================
\section{Summary}
\label{sec:summary}
The acceleration rests on classical linear-algebra identities specialised to the zero-diagonal,
possibly-real matrices produced by maximum-entropy models: $\tr(A^3)/6$ for triangles, a distinct-index
inclusion--exclusion over Hadamard-built dyad matrices for the directed motifs, a marginal-sum Gram identity
for $V$-motifs, a single \texttt{gemv} for ANND, and an exact dihedral orbit reduction for squares (which is
provably not sub-$O(N^4)$ for dense inputs). Every reformulation is proved equivalent to the original for
real matrices, preserves bit-exact integer counts, remains differentiable across all three AD backends, and
was engineered to keep \emph{peak} memory from scaling with the outer thread count. The measured gains range
from stable constant factors ($5$--$22\times$) to growing asymptotic improvements
($192\times$ for ANND and $2482\times$ for $V$-motifs at $N=10^3$), with the pure-square computation made
tractable at $N=5\times10^4$ through the sparse fast path.
\end{document}
